\scanpage{154}
% archive sheet number: 153
\uline{de $i$}. \deleted{\unsure{D'où} \illegible{2}} C'est l'\uuline{Algèbre de Serre}

Supposons $S=\operatorname{Spec}(A)$, $A$ anneau \deleted{\unsure{normal}} \deleted{\inserted{intègre}} \deleted{pour fixer les idées}. Alors on a
\[
A=\unsure{\varinjlim_i|\underset{i}{\underline{\lim}}}A_i
\]
les $A_i$ les sous-anneaux de type fini, \deleted{\unsure{normaux sur} $\mathbb Z$.} normaux \inserted{lorsque $A$ l'est} si on y tient.

[N.B. \deleted{Si} Si $A$ contient un corps premier, si ce dernier est de car $p>0$, les $A_i$ sont des algèbres de type fini sur $\mathbb F_p$. Si $A\supset\mathbb Q$, \unsure{on peut aussi envisager} \unsure{ce seul cas|le seul cas|ce seul corps|en cet cas}, et écrire
\[
A=\varinjlim_\alpha B_\alpha
\]
les $B_\alpha$ les (sous-\inserted{$\mathbb Q$}-algèbres de $A$ qui sont de type fini; \deleted{\unsure{Mais généralement}} \inserted{un élément de} $\unsure{\mathcal V|\mathscr V}((A_i))$ en définit un de $\mathcal V((B_\alpha))$, \unsure{mais n'oublions|mais il s'en faut} que les objets \deleted{\unsure{de}} de $\mathcal V((B_\alpha))$, qu'on arrive \unsure{à|:} fabriquer proviennent en fait de $\mathcal V((A_i))$, cf exemples plus haut)\unsure{,|.} On dira
% slanted
\margin{N.B. Pour la formation des \uline{algèbres de Serre}, le \unsure{passage} des $A_i$ aux $B_\alpha$ est sans effet.}
